\chapter{Keeping Fit}

\section*{Definition and Overview}
Keeping fit means maintaining a level of physical fitness that supports health, daily function, and quality of life. Fitness has several components: aerobic endurance, strength, flexibility, and balance, and a balanced program includes activity for each. The Australian guidelines recommend at least 150 minutes of moderate activity, or 75 minutes of vigorous activity, each week, plus strength exercises on two or more days. Regular physical activity reduces the risk of heart disease, stroke, diabetes, many cancers, depression, and falls, and it helps with weight, sleep, and energy. Fitness can be built at any age, and the key is finding activity that is enjoyable and sustainable.

\section*{Epidemiology}
Around half of Australian adults do not meet the recommended levels of physical activity, and physical inactivity is a leading risk factor for chronic disease and premature death. Rates of inactivity are higher in women, older people, and people with lower incomes, and they rise with age. The health gap between active and inactive people is substantial: meeting the guidelines is associated with a 20 to 30 per cent lower risk of death from all causes. Increasing physical activity even a little brings meaningful benefits, making it one of the highest-value changes a person can make.

\section*{Aetiology and Pathophysiology}
Exercise improves health through a wide range of mechanisms. Aerobic activity strengthens the heart and lungs, improves blood pressure and cholesterol, and improves the body's handling of glucose, reducing the risk of diabetes. Strength training builds muscle, which supports metabolism, protects the joints, and maintains function and independence in older age. Weight-bearing activity strengthens bones and prevents osteoporosis. Flexibility and balance exercises reduce stiffness and falls. Exercise also releases endorphins, improves sleep, and reduces stress, anxiety, and depression. The body adapts to repeated activity, so fitness improves gradually with regular training, and detraining reverses these gains within weeks.

\section*{Clinical Features}
\begin{itemize}
    \item Improved endurance: everyday tasks feel easier
    \item Increased strength and ease with lifting, carrying, and stairs
    \item Better balance, posture, and flexibility
    \item Improved sleep, mood, and energy
    \item Weight management and reduced body fat
    \item Lower blood pressure, cholesterol, and blood glucose
    \item Guidelines: 150 minutes of moderate or 75 minutes of vigorous activity weekly, plus strength training
    \item Fitness is built progressively, and gains continue with consistency
\end{itemize}

\section*{Differential Diagnosis}
Fitness goals and programs are individualised.
\begin{itemize}
    \item \textbf{Aerobic fitness:} walking, running, cycling, and swimming
    \item \textbf{Strength:} resistance training with weights, bands, or body weight
    \item \textbf{Flexibility:} stretching and yoga
    \item \textbf{Balance:} tai chi, standing balance, and stability work
    \item \textbf{Functional fitness:} activities that support daily life and sport
    \item \textbf{Medical considerations:} health conditions affect what is safe and beneficial
\end{itemize}

\section*{When to Seek Medical Care}
Most people can start exercising safely, but those with heart disease, diabetes, or other chronic conditions, and those who have been very inactive, should discuss an exercise plan with their doctor. Anyone who develops chest pain, breathlessness, or other warning symptoms during exercise should stop and seek assessment. Exercise programs for rehabilitation should be guided by health professionals.

\textbf{Call 000 (or local emergency services) for:}
\begin{itemize}
    \item Chest pain, pressure, or tightness during exercise
    \item Collapse, fainting, or confusion
    \item Sudden severe breathlessness or a racing, irregular heartbeat
    \item Severe injury, such as a suspected fracture or head injury
    \item Heat stroke: confusion, very high temperature, or collapse in hot weather
\end{itemize}

\section*{Diagnosis and Investigations}
\begin{itemize}
    \item \textbf{Health assessment:} medical history, blood pressure, and risk factors
    \item \textbf{Fitness assessment:} aerobic capacity, strength, and flexibility
    \item \textbf{Blood tests:} glucose, cholesterol, and other measures as indicated
    \item \textbf{Exercise testing:} for people with heart disease or risk factors
    \item \textbf{Planning:} an individualised program with realistic goals
    \item \textbf{Monitoring:} tracking progress and adjusting the program
\end{itemize}

\section*{Management}
\begin{itemize}
    \item \textbf{Aerobic activity:} brisk walking, jogging, cycling, or swimming most days
    \item \textbf{Strength training:} resistance exercises two or more days a week
    \item \textbf{Flexibility and balance:} stretching, yoga, or tai chi
    \item \textbf{Gradual progression:} building duration and intensity over weeks
    \item \textbf{Variety:} mixing activities to maintain interest and cover all components
    \item \textbf{Social support:} exercising with friends, groups, or a trainer
    \item \textbf{Incorporating activity:} active transport, stairs, and household tasks
    \item \textbf{Recovery:} rest days, sleep, and nutrition
    \item \textbf{Professional guidance:} accredited exercise physiologists and physiotherapists for specific needs
\end{itemize}

\section*{Complications}
\begin{itemize}
    \item Injuries from overuse, poor technique, or excessive progression
    \item Cardiac events during unaccustomed intense exercise in at-risk people
    \item Heat illness and dehydration
    \item Worsening of existing conditions if programs are not tailored
    \item Loss of motivation and stopping activity
    \item The health consequences of inactivity if exercise is not sustained
\end{itemize}

\section*{Prognosis and Outcomes}
The benefits of keeping fit are among the largest in preventive medicine: meeting activity guidelines lowers the risk of premature death by a fifth to a third, and fitness improves mood, sleep, energy, and independence at all ages. Gains come quickly, with measurable improvement in fitness within weeks, and they are sustained with consistency. Even modest increases in activity in previously inactive people produce meaningful benefits.

\section*{Prevention and Self-Care}
\begin{itemize}
    \item Aim for 150 minutes of moderate or 75 minutes of vigorous activity each week
    \item Add strength training on two or more days
    \item Start gradually and build up over weeks
    \item Find activities you enjoy and can sustain
    \item Build activity into daily life: walk, cycle, and take the stairs
    \item Include variety to cover endurance, strength, flexibility, and balance
    \item Exercise with others for motivation
    \item Listen to your body and rest when needed
    \item Get a health check before starting if you have risk factors
    \item Stay hydrated and protect against heat and sun
\end{itemize}

\section*{Sources and Further Reading}
The following reputable sources provide further information for healthcare students and patients seeking reliable, current advice about keeping fit.
\begin{itemize}
    \item Department of Health and Aged Care. \textit{Physical activity guidelines.} https://www.health.gov.au/
    \item Healthdirect Australia. \textit{Physical activity.} https://www.healthdirect.gov.au/
    \item Exercise and Sports Science Australia. \textit{Exercise guidance.} https://www.essa.org.au/
    \item Heart Foundation Australia. \textit{Being active.} https://www.heartfoundation.org.au/
    \item NHS. \textit{Physical activity guidelines.} https://www.nhs.uk/
    \item World Health Organization. \textit{Physical activity.} https://www.who.int/
\end{itemize}
